\subsection{Refinement and coarse-graining channels}

% **Local fix (zero-weight quotient):** the source quotient is completed on
% zero-weight pairs. Documented in
% docs/paper-gaps/cpgsv17_mpdo_zero_weight_preparation_completion.tex.
\begin{theorem}[Positivity and zero-weight neighboring operators]
    \label{thm:mpdo_physical_sector_three_site_neighboring_operator_properties}
    \lean{MPOTensor.PhysicalSectorFactorization.NeighboringTraceFactorization.sector_nonempty,
      MPOTensor.PhysicalSectorFactorization.NeighboringTraceFactorization.threeSiteMiddleIndex_nonempty}
    \lean{MPOTensor.PhysicalSectorFactorization.NeighboringTraceFactorization.threeSiteNeighboringOperator_pos,
      MPOTensor.PhysicalSectorFactorization.NeighboringTraceFactorization.threeSiteNeighboringOperator_eq_zero_of_mul_eq_zero}
    \leanok
    \uses{def:mpdo_physical_sector_three_site_neighboring_operator,
      def:mpdo_neighboring_trace_factorization,
      thm:mpdo_physical_sector_three_site_neighboring_operator_trace}
    The sector set and every middle-subspin space are nonempty. Each
    $\Omega_{k,h}$ is positive semidefinite, and
    $a_kb_h=0$ implies $\Omega_{k,h}=0$.
\end{theorem}

\begin{definition}[The three-site neighboring-state preparation $\mathcal T_1$]
    \label{def:mpdo_physical_sector_t1_preparation}
    \lean{MPOTensor.PhysicalSectorFactorization.NeighboringTraceFactorization.normalizedThreeSiteNeighboringDensity}
    \lean{MPOTensor.PhysicalSectorFactorization.NeighboringTraceFactorization.inactiveThreeSiteNeighboringDensity}
    \lean{MPOTensor.PhysicalSectorFactorization.NeighboringTraceFactorization.completedThreeSiteNeighboringDensity}
    \lean{MPOTensor.PhysicalSectorFactorization.NeighboringTraceFactorization.completedThreeSitePreparationMap}
    \lean{MPOTensor.PhysicalSectorFactorization.NeighboringTraceFactorization.completedThreeSiteControlledMap}
    \lean{MPOTensor.PhysicalSectorFactorization.NeighboringTraceFactorization.t1Map}
    \leanok
    \uses{def:mpdo_physical_sector_three_site_neighboring_operator,
      def:mpdo_neighboring_trace_factorization,
      thm:mpdo_physical_sector_three_site_neighboring_operator_properties}
    On a sector pair with $a_kb_h\ne0$, set
    $\widehat\Omega_{k,h}=(a_kb_h)^{-1}\Omega_{k,h}$.
    If $a_kb_h=0$, positivity and vanishing trace imply
    $\Omega_{k,h}=0$; choose any density operator on that summand and denote
    the resulting completed density by $\overline\Omega_{k,h}$. The sectorwise
    preparation is $X\longmapsto X\otimes\overline\Omega_{k,h}$.
    Taking the orthogonal direct sum over the outer sectors $(k,h)$ gives the
    preparation stage $\mathcal T_1$ of
    \cite[Appendix~C.2, lines~1523--1535]{Cirac2016MPDO_arXiv}.
    The zero-weight choice makes the source map total without imposing a
    nonzero-weight hypothesis.
\end{definition}

\begin{theorem}[Properties of the three-site neighboring densities]
    \label{thm:mpdo_physical_sector_t1_density_properties}
    \lean{MPOTensor.PhysicalSectorFactorization.NeighboringTraceFactorization.normalizedThreeSiteNeighboringDensity_pos,
      MPOTensor.PhysicalSectorFactorization.NeighboringTraceFactorization.trace_normalizedThreeSiteNeighboringDensity}
    \lean{MPOTensor.PhysicalSectorFactorization.NeighboringTraceFactorization.inactiveThreeSiteNeighboringDensity_posDef,
      MPOTensor.PhysicalSectorFactorization.NeighboringTraceFactorization.inactiveThreeSiteNeighboringDensity_trace}
    \lean{MPOTensor.PhysicalSectorFactorization.NeighboringTraceFactorization.completedThreeSiteNeighboringDensity_pos,
      MPOTensor.PhysicalSectorFactorization.NeighboringTraceFactorization.trace_completedThreeSiteNeighboringDensity,
      MPOTensor.PhysicalSectorFactorization.NeighboringTraceFactorization.completedThreeSiteNeighboringDensity_eq_normalized}
    \lean{MPOTensor.PhysicalSectorFactorization.NeighboringTraceFactorization.completedThreeSitePreparationMap_eq_normalized}
    \lean{MPOTensor.PhysicalSectorFactorization.NeighboringTraceFactorization.completedThreeSiteControlledMap_sameBlock_apply}
    \leanok
    \uses{def:mpdo_physical_sector_t1_preparation}
    The normalized density is positive semidefinite and has trace one on
    every active pair. The density chosen on an inactive pair is positive
    definite and has trace one. Hence $\overline\Omega_{k,h}$ is positive
    semidefinite with trace one for every pair, and agrees with
    $\widehat\Omega_{k,h}$ on active pairs. The corresponding sectorwise and
    controlled maps prepare these densities on their diagonal summands.
\end{theorem}

\begin{theorem}[$\mathcal T_1$ is trace-preserving completely positive]
    \label{thm:mpdo_physical_sector_t1_preparation_cptp}
    \lean{MPOTensor.PhysicalSectorFactorization.NeighboringTraceFactorization.completedThreeSitePreparationMap_isKrausCPTP}
    \lean{MPOTensor.PhysicalSectorFactorization.NeighboringTraceFactorization.completedThreeSiteControlledMap_isKrausCPTP}
    \lean{MPOTensor.PhysicalSectorFactorization.NeighboringTraceFactorization.t1Map_isKrausCPTP}
    \leanok
    \uses{def:mpdo_physical_sector_t1_preparation,
      thm:mpdo_physical_sector_t1_density_properties}
    Each completed preparation is trace-preserving and completely positive.
    The same holds for their orthogonally controlled direct sum.
\end{theorem}

\begin{proof}
    \leanok
    On an active pair,
    $\tr(\widehat\Omega_{k,h})=(a_kb_h)^{-1}a_kb_h=1$; the chosen density on
    a zero-weight pair also has trace one. Preparation by a positive operator
    of trace one is trace-preserving and completely positive. The orthogonal
    control discards coherences between distinct outer-sector pairs and
    applies these preparations on the diagonal summands.
\end{proof}

\begin{definition}[The refinement channel $\mathcal T$]
    \label{def:mpdo_physical_sector_t_map}
    \lean{MPOTensor.PhysicalSectorFactorization.NeighboringTraceFactorization.tMap}
    \leanok
    \uses{def:mpdo_t0_map,
      def:mpdo_physical_sector_t1_preparation,
      def:mpdo_t2_map}
    Define the channel from two physical sites to three physical sites by
    $\mathcal T=\mathcal T_2\mathcal T_1\mathcal T_0$.
    Thus $\mathcal T_0$ traces the neighboring subspins, $\mathcal T_1$
    adjoins the completed three-site neighboring density, and $\mathcal T_2$
    restores the three complete physical sectors. This is the refinement map
    in \cite[Appendix~C.2, lines~1522--1545]{Cirac2016MPDO_arXiv}.
\end{definition}

\begin{theorem}[$\mathcal T$ is trace-preserving completely positive]
    \label{thm:mpdo_physical_sector_t_map_cptp}
    \lean{MPOTensor.PhysicalSectorFactorization.NeighboringTraceFactorization.tMap_isKrausCPTP}
    \leanok
    \uses{def:mpdo_physical_sector_t_map}
    The refinement channel $\mathcal T$ is trace-preserving and completely
    positive.
\end{theorem}

\begin{proof}
    \leanok
    \uses{thm:mpdo_t0_map_cptp,
      thm:mpdo_physical_sector_t1_preparation_cptp,
      thm:mpdo_t2_map_cptp,
      lem:is_kraus_cptp_comp}
    Each constituent map is trace-preserving and completely positive, and
    this property is preserved under composition.
\end{proof}

Applied to the two-site closure, on the active $(k,h)$-block,
$\mathcal T_0$ traces $B_k^R\otimes B_h^L$ and sends
$\mathfrak R_2(\mathcal K_2(X))_{k,h}$ to
$a_kb_hB_{k,h}(X)$ on $B_k^L\otimes B_h^R$. The preparation
$\mathcal T_1$ then adjoins
\begin{align}
  \widehat\Omega_{k,h}
  &=(a_kb_h)^{-1}\bigoplus_l
    (\eta_{k,l}\otimes\eta_{l,h}),
  \notag
\end{align}
and $\mathcal T_2$ orders the result as
\begin{align}
  (B_k^L\otimes B_k^R)\otimes(B_l^L\otimes B_l^R)
  \otimes(B_h^L\otimes B_h^R).
  \notag
\end{align}
The refinement uses the sitewise sector coordinates $\mathfrak R_2$ and
$\mathfrak R_3$
\cite[Appendix~C.2, lines~1510--1516 and~1522--1545]
  {Cirac2016MPDO_arXiv}.
For each outer pair, write
\begin{align}
  \mathcal H^{\Omega}_{k,h}
  &=\bigoplus_l
    (B_k^R\otimes B_l^L)\otimes(B_l^R\otimes B_h^L)
  \notag
\end{align}
for the carrier Hilbert space of $\Omega_{k,h}$. The symbols
$\widehat\Omega_{k,h}$ and $\overline\Omega_{k,h}$ denote respectively its
normalized active density and its completed density on every sector pair.
With the matrix algebras introduced above, the refinement stages have type

% Formula: T_0 : A_2 -> A_partial, T_1 : A_partial -> A_Omega, and
% T_2 : A_Omega -> A_3. Ink: each blue stroke denotes a channel, its label
% names the stage, and the fixed left-to-right axis records composition.
\par\noindent\hfil
  \begin{tikzcd}[channel, column sep=large]
    \mathcal A_2 \arrow[r, "\mathcal T_0"] &
    \mathcal A_{\partial} \arrow[r, "\mathcal T_1"] &
    \mathcal A_{\Omega} \arrow[r, "\mathcal T_2"] &
    \mathcal A_3
  \end{tikzcd}
\hfil\par
Here $\mathcal T_1$ adjoins $\overline\Omega_{k,h}$ on each diagonal
outer-sector summand. The composite channel has type
% Formula: T = T_2 T_1 T_0 : A_2 -> A_3. Ink: the blue stroke denotes the
% composite channel T; its endpoints are matrix algebras, not matrices.
\par\noindent\hfil
  \begin{tikzcd}[channel, column sep=large]
    \mathcal A_2 \arrow[r, "\mathcal T"] & \mathcal A_3
  \end{tikzcd}
\hfil\par

\begin{definition}[The neighboring-state preparation $\mathcal S_1$]
    \label{def:mpdo_physical_sector_s1_map}
    \lean{MPOTensor.PhysicalSectorFactorization.NeighboringTraceFactorization.normalizedNeighboringDensity}
    \lean{MPOTensor.PhysicalSectorFactorization.NeighboringTraceFactorization.normalizedNeighboringDensity_pos,
      MPOTensor.PhysicalSectorFactorization.NeighboringTraceFactorization.trace_normalizedNeighboringDensity,
      MPOTensor.PhysicalSectorFactorization.NeighboringTraceFactorization.normalizedNeighboringPreparationMap_isKrausCPTP}
    \lean{MPOTensor.PhysicalSectorFactorization.NeighboringTraceFactorization.inactiveNeighboringDensity,
      MPOTensor.PhysicalSectorFactorization.NeighboringTraceFactorization.inactiveNeighboringDensity_posDef,
      MPOTensor.PhysicalSectorFactorization.NeighboringTraceFactorization.inactiveNeighboringDensity_trace}
    \lean{MPOTensor.PhysicalSectorFactorization.NeighboringTraceFactorization.completedNeighboringDensity}
    \lean{MPOTensor.PhysicalSectorFactorization.NeighboringTraceFactorization.completedNeighboringDensity_pos,
      MPOTensor.PhysicalSectorFactorization.NeighboringTraceFactorization.trace_completedNeighboringDensity,
      MPOTensor.PhysicalSectorFactorization.NeighboringTraceFactorization.completedNeighboringDensity_eq_normalized}
    \lean{MPOTensor.PhysicalSectorFactorization.NeighboringTraceFactorization.completedNeighboringPreparationMap,
      MPOTensor.PhysicalSectorFactorization.NeighboringTraceFactorization.completedNeighboringPreparationMap_isKrausCPTP,
      MPOTensor.PhysicalSectorFactorization.NeighboringTraceFactorization.completedNeighboringPreparationMap_eq_normalized}
    \lean{MPOTensor.PhysicalSectorFactorization.NeighboringTraceFactorization.completedNeighboringControlledMap}
    \lean{MPOTensor.PhysicalSectorFactorization.NeighboringTraceFactorization.completedNeighboringControlledMap_sameBlock_apply}
    \lean{MPOTensor.PhysicalSectorFactorization.NeighboringTraceFactorization.completedNeighboringPreparationMap_smul}
    \lean{MPOTensor.PhysicalSectorFactorization.NeighboringTraceFactorization.s1Map}
    \lean{MPOTensor.PhysicalSectorFactorization.NeighboringTraceFactorization.s1Map_block_eq_kronecker,
      MPOTensor.PhysicalSectorFactorization.NeighboringTraceFactorization.s1Map_apply_of_ne}
    \leanok
    \uses{def:mpdo_neighboring_trace_factorization}
    On a sector pair with $a_kb_h\ne0$, let
    $\mathcal S_{k,h}(X)=X\otimes\frac{\eta_{k,h}}{a_kb_h}$.
    If $a_kb_h=0$, positivity and vanishing trace imply
    $\eta_{k,h}=0$; choose any density operator on that unused summand.
    Taking the orthogonal direct sum over $(k,h)$ defines $\mathcal S_1$.
\end{definition}

\begin{theorem}[$\mathcal S_1$ is trace-preserving completely positive]
    \label{thm:mpdo_physical_sector_s1_map_cptp}
    \lean{MPOTensor.PhysicalSectorFactorization.NeighboringTraceFactorization.completedNeighboringControlledMap_isKrausCPTP}
    \lean{MPOTensor.PhysicalSectorFactorization.NeighboringTraceFactorization.s1Map_isKrausCPTP}
    \leanok
    \uses{def:mpdo_physical_sector_s1_map}
    The map $\mathcal S_1$ is trace-preserving and completely positive.
\end{theorem}

\begin{definition}[The coarse-graining channel $\mathcal S$]
    \label{def:mpdo_physical_sector_s_map}
    \lean{MPOTensor.PhysicalSectorFactorization.NeighboringTraceFactorization.sMap}
    \leanok
    \uses{def:mpdo_s0_map, def:mpdo_physical_sector_s1_map,
      def:mpdo_s2_map}
    Define the channel from three physical sites to two physical sites by
    $\mathcal S=\mathcal S_2\mathcal S_1\mathcal S_0$.
    This is the coarse-graining map in Proposition~C.7
    \cite[Appendix~C.2, lines~1547--1563]{Cirac2016MPDO_arXiv}.
\end{definition}

\begin{theorem}[$\mathcal S$ is trace-preserving completely positive]
    \label{thm:mpdo_physical_sector_s_map_cptp}
    \lean{MPOTensor.PhysicalSectorFactorization.NeighboringTraceFactorization.sMap_isKrausCPTP}
    \leanok
    \uses{thm:mpdo_s0_map_cptp,
      thm:mpdo_physical_sector_s1_map_cptp,
      thm:mpdo_s2_map_cptp,
      lem:is_kraus_cptp_comp}
    The channel $\mathcal S$ is trace-preserving and completely positive.
\end{theorem}

\begin{proof}
    \leanok
    \uses{lem:is_kraus_cptp_comp}
    Each of the three factors is trace-preserving and completely positive,
    and this property is preserved under composition.
\end{proof}

\begin{definition}[Two- and three-site sector coordinates]
    \label{def:mpdo_physical_sector_closure_coordinates}
    \lean{MPOTensor.PhysicalSectorFactorization.twoSiteSectorCoordinateEquiv,
      MPOTensor.PhysicalSectorFactorization.twoSiteSectorCoordinateEquiv_symm_apply,
      MPOTensor.PhysicalSectorFactorization.threeSiteSectorCoordinateEquiv,
      MPOTensor.PhysicalSectorFactorization.threeSiteSectorCoordinateEquiv_symm_apply}
    \leanok
    \uses{def:mpdo_two_site_sector_closure,
      def:mpdo_three_site_sector_closure}
    The coordinate equivalences used by the coarse-graining channel expose
    the direct sums over site-sector pairs and triples. Write
    $\mathfrak R_2$ and $\mathfrak R_3$ for the corresponding matrix
    reindexings. Their inverses send every sector block to the corresponding
    physical matrix entries. These reindexings preserve the left-right order
    within every physical site. They differ from the regrouped coordinates
    $\mathfrak G_2$ and $\mathfrak G_3$ of
    Definition~\ref{def:mpdo_global_sector_closure_coordinates}, which place
    the two outer subspins before the neighboring subspins. If
    $R_{\Phi_2}$ and $R_{\Phi_3}$ denote the two regrouping maps, then
    \begin{align}
      \mathfrak G_2&=R_{\Phi_2}\mathfrak R_2,
      \notag\\
      \mathfrak G_3&=R_{\Phi_3}\mathfrak R_3.
      \notag
    \end{align}
\end{definition}

\begin{theorem}[Diagonal outer-sector block of the three-site closure]
    \label{thm:mpdo_regrouped_three_site_closure_block}
    \lean{MPOTensor.PhysicalSectorFactorization.s0Regrouped_threeSiteClosure_block_eq}
    \leanok
    \uses{def:phys_close,
      def:mpdo_sector_boundary_operator,
      def:mpdo_physical_sector_three_site_neighboring_operator,
      def:mpdo_physical_sector_closure_coordinates,
      def:mpdo_three_site_subspin_regrouping}
    For every virtual matrix $X$ and outer-sector pair $(k,h)$,
    \begin{align}
      \left[R_{\Phi_3}\mathfrak R_3(\mathcal K_3(X))\right]_{(k,h),(k,h)}
      &=B_{k,h}(X)\otimes\Omega_{k,h}.
      \notag
    \end{align}
\end{theorem}

\begin{proof}
    \leanok
    \uses{thm:mpdo_three_site_sector_closure_factorization,
      def:mpdo_physical_sector_factorization}
    On the summand labelled by the middle sector $l$, this is the fixed-sector
    factorization
    $B_{k,h}(X)\otimes(\eta_{k,l}\otimes\eta_{l,h})$.
    Taking the direct sum over $l$ gives the stated block.
\end{proof}

\begin{theorem}[Off-diagonal outer-sector entries of the three-site closure]
    \label{thm:mpdo_regrouped_three_site_closure_off_diagonal}
    \lean{MPOTensor.PhysicalSectorFactorization.s0Regrouped_threeSiteClosure_apply_of_outer_ne}
    \leanok
    \uses{def:mpdo_physical_sector_factorization,
      def:mpdo_physical_sector_closure_coordinates,
      def:mpdo_three_site_subspin_regrouping}
    Every matrix entry of
    $R_{\Phi_3}\mathfrak R_3(\mathcal K_3(X))$ between distinct outer-sector
    pairs $(k,h)\ne(p,q)$ vanishes.
\end{theorem}

\begin{proof}
    \leanok
    \uses{def:mpdo_physical_sector_factorization}
    Distinct outer-sector pairs differ either in their first sector or in
    their third sector. In the corresponding physical slice, the
    sector-coordinate tensor is block diagonal, so that factor vanishes.
\end{proof}

\begin{theorem}[Regrouped two-site closure block]
    \label{thm:mpdo_t0_regrouped_two_site_closure_block}
    \lean{MPOTensor.PhysicalSectorFactorization.t0Regrouped_twoSiteClosure_block_eq}
    \leanok
    \uses{thm:mpdo_two_site_sector_closure_factorization,
      def:mpdo_physical_sector_closure_coordinates,
      def:mpdo_two_site_subspin_regrouping}
    For every virtual matrix $X$ and outer-sector pair $(k,h)$,
    \begin{align}
      \left[R_{\Phi_2}\mathfrak R_2(\mathcal K_2(X))\right]_{(k,h),(k,h)}
      &=B_{k,h}(X)\otimes\eta_{k,h}.
      \notag
    \end{align}
\end{theorem}

\begin{proof}
    \leanok
    \uses{thm:mpdo_two_site_sector_closure_factorization}
    This is the fixed-sector two-site factorization after identifying the
    sitewise sector coordinates with the regrouped boundary and neighboring
    coordinates.
\end{proof}

\begin{theorem}[Diagonal-block action of $\mathcal T_0$]
    \label{thm:mpdo_t0_same_block_action}
    \lean{MPOTensor.PhysicalSectorFactorization.t0Map_sameBlock_apply}
    \leanok
    \uses{def:mpdo_t0_map,
      def:mpdo_two_site_subspin_regrouping}
    For every two-site matrix $Y$, outer-sector pair $(k,h)$, and boundary
    indices $a,b$,
    \begin{align}
      [\mathcal T_0(Y)]_{(k,h,a),(k,h,b)}
      &=
      \left[\tr_{R_k\otimes L_h}
          \left([R_{\Phi_2}Y]_{(k,h),(k,h)}\right)\right]_{a,b}.
      \notag
    \end{align}
\end{theorem}

\begin{proof}
    \leanok
    The controlled partial trace acts on a diagonal outer-sector block by
    the ordinary partial trace over its neighboring factor.
\end{proof}

\begin{theorem}[Partial trace of a factorized two-site block]
    \label{thm:mpdo_t0_factorized_block_action}
    \lean{MPOTensor.PhysicalSectorFactorization.NeighboringTraceFactorization.t0Map_block_eq_smul}
    \leanok
    \uses{thm:mpdo_t0_same_block_action,
      def:mpdo_neighboring_trace_factorization}
    If the $(k,h)$ diagonal block of $R_{\Phi_2}Y$ is
    $B\otimes\eta_{k,h}$, then
    $[\mathcal T_0(Y)]_{(k,h),(k,h)}=a_kb_hB$.
\end{theorem}

\begin{proof}
    \leanok
    \uses{thm:mpdo_t0_same_block_action,
      def:mpdo_neighboring_trace_factorization}
    The partial trace of $B\otimes\eta_{k,h}$ is
    $\tr(\eta_{k,h})B=a_kb_hB$.
\end{proof}

% **Local fix (zero-weight quotient):** the next three entries use the
% completion documented in
% docs/paper-gaps/cpgsv17_mpdo_zero_weight_preparation_completion.tex.
\begin{theorem}[Preparation of a scaled boundary block]
    \label{thm:mpdo_t1_scaled_block_preparation}
    \lean{MPOTensor.PhysicalSectorFactorization.NeighboringTraceFactorization.completedThreeSitePreparationMap_smul}
    \leanok
    \uses{def:mpdo_physical_sector_t1_preparation,
      thm:mpdo_physical_sector_t1_density_properties}
    For every outer-sector pair $(k,h)$ and boundary matrix $B$, the completed
    sectorwise preparation satisfies
    $\overline{\mathcal T}_{k,h}(a_kb_hB)
    =B\otimes\Omega_{k,h}$.
\end{theorem}

\begin{proof}
    \leanok
    \uses{thm:mpdo_physical_sector_t1_density_properties,
      thm:mpdo_physical_sector_three_site_neighboring_operator_trace}
    If $a_kb_h\ne0$, the scalar cancels the normalization in
    $\widehat\Omega_{k,h}=(a_kb_h)^{-1}\Omega_{k,h}$. If $a_kb_h=0$, both
    sides vanish because $\Omega_{k,h}=0$.
\end{proof}

\begin{theorem}[Diagonal-block action of $\mathcal T_1$]
    \label{thm:mpdo_t1_factorized_block_action}
    \lean{MPOTensor.PhysicalSectorFactorization.NeighboringTraceFactorization.t1Map_block_eq_kronecker}
    \leanok
    \uses{def:mpdo_physical_sector_t1_preparation,
      thm:mpdo_t1_scaled_block_preparation}
    If the $(k,h)$ diagonal block of a boundary matrix $Y$ is $a_kb_hB$, then
    $[\mathcal T_1(Y)]_{(k,h),(k,h)}
    =B\otimes\Omega_{k,h}$.
\end{theorem}

\begin{proof}
    \leanok
    \uses{thm:mpdo_t1_scaled_block_preparation}
    On a diagonal outer-sector block, the controlled map is the completed
    sectorwise preparation, to which the preceding theorem applies.
\end{proof}

\begin{theorem}[Off-diagonal action of $\mathcal T_1$]
    \label{thm:mpdo_t1_off_diagonal_action}
    \lean{MPOTensor.PhysicalSectorFactorization.NeighboringTraceFactorization.t1Map_apply_of_ne}
    \leanok
    \uses{def:mpdo_physical_sector_t1_preparation}
    The map $\mathcal T_1$ annihilates every matrix entry between distinct
    outer-sector pairs.
\end{theorem}

\begin{proof}
    \leanok
    This is the orthogonal control in the definition of $\mathcal T_1$.
\end{proof}

\begin{theorem}[Reindexing action of $\mathcal T_2$]
    \label{thm:mpdo_t2_reindexing_action}
    \lean{MPOTensor.PhysicalSectorFactorization.t2Map_apply}
    \leanok
    \uses{def:mpdo_t2_map}
    Let $\Phi_3$ be the three-site regrouping. For every prepared matrix $Y$
    and three-site indices $p,q$,
    $[\mathcal T_2(Y)]_{p,q}=Y_{\Phi_3(p),\Phi_3(q)}$.
\end{theorem}

\begin{proof}
    \leanok
    This is matrix reindexing along the inverse regrouping
    $\Phi_3^{-1}$.
\end{proof}

\begin{theorem}[Refinement of the two-site closure]
    \label{thm:mpdo_physical_sector_t_closure_identity}
    \lean{MPOTensor.PhysicalSectorFactorization.NeighboringTraceFactorization.tMap_twoSiteClosure_eq_threeSiteClosure}
    \leanok
    \uses{def:phys_close,
      def:mpdo_physical_sector_t_map,
      def:mpdo_physical_sector_closure_coordinates,
      def:mpdo_neighboring_trace_factorization}
    Assume the neighboring-operator trace factorization of
    Definition~\ref{def:mpdo_neighboring_trace_factorization}. For every
    virtual matrix $X$, in the canonical physical-sector coordinates selected
    by the isometry $U$, the refinement channel satisfies
    \begin{align}
      \mathcal T(\mathfrak R_2(\mathcal K_2(X)))
      &=\mathfrak R_3(\mathcal K_3(X)).
      \label{eq:rfp_t_closure}
    \end{align}
    This is the refinement half of Proposition~C.7
    \cite[Appendix~C.2, lines~1510--1516 and~1522--1545]
      {Cirac2016MPDO_arXiv}.
\end{theorem}

\begin{proof}
    \leanok
    \uses{thm:mpdo_t0_regrouped_two_site_closure_block,
      thm:mpdo_t0_factorized_block_action,
      thm:mpdo_t1_factorized_block_action,
      thm:mpdo_t1_off_diagonal_action,
      thm:mpdo_regrouped_three_site_closure_block,
      thm:mpdo_regrouped_three_site_closure_off_diagonal}
    The regrouping contained in $\mathcal T_0$ changes
    $\mathfrak R_2$ into $\mathfrak G_2$. On the $(k,h)$ diagonal block, the
    partial trace then gives
    \begin{align}
      \mathcal T_0
        (\mathfrak R_2(\mathcal K_2(X)))_{k,h}
      &=a_kb_hB_{k,h}(X).
      \notag
    \end{align}
    On an active pair, $\mathcal T_1$ adjoins
    $(a_kb_h)^{-1}\Omega_{k,h}$, and hence
    \begin{align}
      \mathcal T_1\mathcal T_0
        (\mathfrak R_2(\mathcal K_2(X)))_{k,h}
      &=B_{k,h}(X)\otimes\Omega_{k,h}.
      \notag
    \end{align}
    If $a_kb_h=0$, both sides vanish because
    $\Omega_{k,h}=0$. After the canonical reassociation of the intermediate
    sector sum, these are precisely the diagonal blocks of
    $\mathfrak G_3(\mathcal K_3(X))$; all off-diagonal sector blocks vanish.
    Finally, the inverse regrouping $\mathcal T_2$ changes
    $\mathfrak G_3$ into $\mathfrak R_3$, which gives
    (\ref{eq:rfp_t_closure}).
\end{proof}

\begin{theorem}[Coarse-graining of the three-site closure]
    \label{thm:mpdo_physical_sector_s_closure_identity}
    \lean{MPOTensor.PhysicalSectorFactorization.NeighboringTraceFactorization.s0Map_block_eq_smul}
    \lean{MPOTensor.PhysicalSectorFactorization.NeighboringTraceFactorization.sMap_threeSiteClosure_eq_twoSiteClosure}
    \leanok
    \uses{def:phys_close, def:mpdo_physical_sector_s_map,
      def:mpdo_physical_sector_closure_coordinates,
      def:mpdo_neighboring_trace_factorization}
    Assume the neighboring-operator trace factorization of
    Definition~\ref{def:mpdo_neighboring_trace_factorization}. For every
    virtual matrix $X$, in the canonical physical-sector
    coordinates selected by the isometry $U$, the coarse-graining channel
    satisfies
    \begin{align}
      \mathcal S(\mathfrak R_3(\mathcal K_3(X)))
      &=\mathfrak R_2(\mathcal K_2(X)).
      \label{eq:rfp_s_closure}
    \end{align}
\end{theorem}

\begin{proof}
    \leanok
    \uses{thm:mpdo_regrouped_three_site_closure_block,
      thm:mpdo_two_site_sector_closure_factorization,
      thm:mpdo_three_site_sector_closure_factorization,
      thm:mpdo_physical_sector_three_site_neighboring_operator_trace,
      def:mpdo_physical_sector_factorization,
      def:mpdo_physical_sector_closure_coordinates,
      def:mpdo_physical_sector_s1_map}
    On the $(k,h)$ diagonal block, $\mathcal S_0$ traces
    $\Omega_{k,h}$ and gives $a_kb_hB_{k,h}(X)$. On an active sector pair,
    $\mathcal S_1((a_kb_h)B_{k,h}(X))
    =B_{k,h}(X)\otimes\eta_{k,h}$.
    If $a_kb_h=0$, positivity and zero trace give $\eta_{k,h}=0$, so both
    sides of this identity vanish. Finally, $\mathcal S_2$ regroups the four
    factors into two physical sites. For distinct sector pairs
    $(k,h)\ne(p,q)$,
    \begin{align}
      (\mathcal S_1Y)_{(k,h),(p,q)}&=0,
      \notag\\
      [\mathfrak R_2(\mathcal K_2(X))]_{(k,h),(p,q)}&=0.
      \notag
    \end{align}
    These diagonal and off-diagonal identities give
    (\ref{eq:rfp_s_closure}).
\end{proof}

For the following construction, through
Theorem~\ref{thm:mpdo_blocked_original_tensor_is_rfp_via_ts}, assume the
neighboring-operator trace factorization of
Definition~\ref{def:mpdo_neighboring_trace_factorization}. Write
$q=\dim\!\left(\bigoplus_k L_k\otimes R_k\right)$
for the physical dimension of the sector-coordinate tensor
$\widehat{\mathcal K}$.
